\section{Finite certificates for failure to be a price}
\label{sec:duality}

This section makes the reduction to wire identities exact. Its
algebraic input is the cover form of Strassen duality proved in
\cite[Theorem A.6, pp.~20--21]{Anthropic2026}; the extension of
ordinary spectral points from subsemirings is also a standard
consequence of Strassen's theorem \cite[Theorem 4.16]{wz-spectra}.
The distinction matters: products of wires alone do not form a
semiring, because they are not closed under direct sum.

\begin{theorem}[Cover duality, after Strassen]\label{thm:coverduality}
Let $C\ne0$ and $A_1,\ldots,A_r$ be tensors on a fixed finite
party set, with at least one $A_j\ne0$.
The following are equivalent:
\begin{enumerate}
\item $\phi(C)\le\max_j\phi(A_j)$ for every spectral point in
Definition~\ref{def:price};
\item for every $\eps>0$ and all sufficiently large integers $N$,
$C^N$ has a cover by at most $2^{\eps N}$ tensors, each a product
of exactly $N$ factors chosen from $A_1,\ldots,A_r$.
\end{enumerate}
\end{theorem}

\begin{proof}
We include the proof because the precise kind of cover is essential.
An upper functional is a nonnegative function $p$ with $p(1)=1$,
submultiplicativity, and subadditivity on covers.
Say it respects $(x,y,c)$ when $p(tx)\le c\,p(ty)$ for every $t$.
The following argument is \cite[Proposition A.5]{Anthropic2026}.

Among upper functionals bounded by a given one and respecting
its triples, choose a pointwise minimal one by Zorn's lemma.
A decreasing chain has its pointwise infimum in this class:
every condition uses only finitely many values, and a common
member of the chain approximates their infima simultaneously.
Write this minimal element as $p$ and put
$r(t)=\lim_n p(t^n)^{1/n}=\inf_n p(t^n)^{1/n}$.
The limit follows from submultiplicativity; if a value is zero,
the same conclusion holds directly.
It satisfies $r(1)=1$ and is submultiplicative.
For a cover of $s$ by $t_1,\ldots,t_l$, expand its $n$th power
into words. For every $\delta>0$ a constant $K_\delta$ satisfies
$p(t_j^a)\le K_\delta(r(t_j)+\delta)^a$ for all $j,a$.
Grouping each word by its multiplicities bounds the total by
$K_\delta^l(\sum_j(r(t_j)+\delta))^n$.
Taking roots and then $\delta\downarrow0$ proves cover
subadditivity for $r$. Iterating a respected triple gives
$p(t^n x^n)\le c^n p(t^n y^n)$, so $r$ respects it.
Minimality implies $r=p$.

For $p(f)>0$, the decreasing sequence
$p(tf^n)/p(f)^n$ has infimum $p_f(t)$.
It defines an upper functional: normalization uses
$p(f^n)=p(f)^n$, and submultiplicativity follows by splitting
the exponent $n$ between the two factors before taking limits.
All respected triples survive the limit. Minimality gives
$p_f=p$, hence $p(tf)=p(t)p(f)$.
If $p(f)=0$, submultiplicativity gives the same equality.
Thus the minimal upper functional is spectral.

Now assume the first assertion and fix $\eps>0$. Set
$\eta=2^{-\eps/2}$ and $A^v=\prod_j A_j^{v_j}$.
A level-$k$ certificate for $t$ is a cover of $tC^k$
by terms $C^{u_i}A^{v_i}$ with $u_i+|v_i|=k$.
Give it weight $\sum_i\eta^{|v_i|}$, and let $p(t)$ be
the infimum of its certificate weights over all levels.
Every tensor has a level-zero certificate by rank-one terms.
Certificates multiply. They can also be raised to a common
level by multiplication by $C$ and combined along covers.
Consequently $p$ is submultiplicative and cover-subadditive,
and $p(1)\le1$.

If $p(1)=1$, it is an upper functional and respects
$(A_j,C,\eta)$: multiply a certificate for $tC$ by $A_j$.
The preceding minimality argument gives a spectral $\phi$
with $\phi(A_j)\le\eta\phi(C)$ for all $j$.
This contradicts the first assertion, since $\phi(C)\ge1$.
Thus $p(1)<1$. Choose a certificate of $1$ of weight less
than $1$ and level $k$. Necessarily $k\ge1$, and every term
contains at least one $A$ factor; otherwise its weight alone
would be $1$.

In $C^N$, repeatedly replace any $k$ remaining $C$ factors
by this certificate. Every replacement strictly reduces the
number of $C$ factors along a branch, and its total weight
does not increase. Each final term has fewer than $k$ copies
of $C$, has $N$ factors altogether, and has weight at least
$\eta^N$. There are therefore at most $\eta^{-N}$ terms.
Fix a nonzero $A_j$. Since $1$ is a restriction of $A_j$,
a fixed number $c$ of copies of $A_j$ covers $C$.
Replace the fewer than $k$ remaining $C$ factors by these
covers. The total is at most $c^k\eta^{-N}\le2^{\eps N}$
for all large $N$, and each term has exactly $N$ $A$ factors.
This proves the second assertion.

Conversely apply any spectral $\phi$ to these covers.
Taking $N$th roots gives
$\phi(C)\le2^\eps\max_j\phi(A_j)$.
Let $\eps\downarrow0$.
\end{proof}

For a finite set $V$ write $E=\binom V2$ and, for
$u\in\N^E$, set
$G_u=\prod_{\{i,j\}\in E}\langle2^{u_{ij}}\rangle_{ij}$.
For a spanning tree $T$ on $V$ let
$\tau_T(d)=\sum_{\{i,j\}\in E(T)}d_{ij}$.

\begin{lemma}[Separation of a closed cone]\label{lem:coneseparation}
Let $K\subseteq\R_{\ge0}^q$ be a closed cone, containing zero,
and let $e\in\R_{\ge0}^q\setminus K$.
There are $u,v_1,\ldots,v_r\in\N^q$, with $r\le2q$, such that
\[
 \langle u,d\rangle\le\max_j\langle v_j,d\rangle\quad(d\in K),
 \qquad
 \langle u,e\rangle>\max_j\langle v_j,e\rangle.
\]
Convexity of $K$ is not required.
\end{lemma}

\begin{proof}
Normalize by $\sigma(x)=\sum_i x_i$.
The intersection of $K$ with the simplex $\sigma(x)=1$
is compact and omits $e/\sigma(e)$.
Choose a box with rational endpoints, containing this point
strictly, whose closure in the simplex avoids that intersection.
If the bounds are $a_i<x_i<b_i$, its homogenized inequalities
are $x_i-a_i\sigma(x)>0$ and $b_i\sigma(x)-x_i>0$.
After a common positive integer multiplication, these are
integer linear forms $\ell_1,\ldots,\ell_r$, all positive at $e$.
For every nonzero $d\in K$, at least one $\ell_j(d)\le0$;
the same weak conclusion holds at zero.
Choose a nonnegative integer vector $u$ whose coordinates
dominate the coefficients of every $\ell_j$ and put
$v_j=u-\ell_j$, regarding a linear form as its coefficient vector.
Then every $v_j$ is nonnegative and the two claims follow.
\end{proof}

\begin{theorem}[A finite certificate for every nonprice]
\label{thm:finitecertificate}
Let $|V|\ge2$, let $\kk$ be a field, and let $e$ be a nonnegative
function on the unordered pairs of $V$. Then $e$ is not a price
over $\kk$ if and only if there exist loads
$u,v_1,\ldots,v_r\in\N^E$, an integer $N\ge1$, a spanning tree
$T$ on $V$, and a cover of $G_u^N$ by $M\ge1$ tensors, each a
product of $N$ factors from $\{G_{v_1},\ldots,G_{v_r}\}$, such that
\begin{equation}\label{eq:finitecertificate}
 N\bigl(\langle u,e\rangle-\max_j\langle v_j,e\rangle\bigr)
   >\tau_T(e)\log_2 M .
\end{equation}
Every finite nonprice is therefore excluded by one finite
algebraic tensor identity with a strict length margin.
\end{theorem}

\begin{proof}
The set of prices is a closed cone by definition.
If $e$ is outside it, Lemma~\ref{lem:coneseparation} supplies
$u,v_j$ and a positive gap
$\delta=\langle u,e\rangle-\max_j\langle v_j,e\rangle$.
For every spectral point, exponentiation of the valid
homogeneous inequality gives
$\phi(G_u)\le\max_j\phi(G_{v_j})$.
Apply Theorem~\ref{thm:coverduality} with any
$\eps>0$ satisfying $\tau_T(e)\eps<\delta$.
For all large $N$ there is the stated cover with
$\log_2M\le\eps N$, proving \eqref{eq:finitecertificate}.
This choice also works when $\tau_T(e)=0$.

Conversely, the terminal-selector inequality, with terminal set
$V$, bounds $\log_2\phi(G_u^N)$ by the largest logarithmic
value of a covering tensor plus $g_\phi(V)\log_2M$.
Lemma~\ref{lem:wire} bounds $g_\phi(V)$ by $\tau_T(d_\phi)$.
Every covering tensor has load equal to a sum of $N$ of the
$v_j$. We obtain
\[
 N\langle u,d_\phi\rangle
 \le N\max_j\langle v_j,d_\phi\rangle+
          \tau_T(d_\phi)\log_2M .
\]
This inequality is homogeneous and continuous, hence holds
for every price. It contradicts \eqref{eq:finitecertificate}.
\end{proof}

\begin{corollary}\label{cor:allpriceineq}
The price cone is the intersection of all homogeneous inequalities
$\langle u,d\rangle\le\max_j\langle v_j,d\rangle$
that arise from the subexponential covers of
Theorem~\ref{thm:coverduality}, with nonnegative integer wire loads.
In particular these inequalities are complete for detecting
failure to be a price, even when a single linear inequality is not.
\end{corollary}

\begin{proof}
Every such inequality holds for prices by
Lemma~\ref{lem:homogeneous}. Lemma~\ref{lem:coneseparation} and
Theorem~\ref{thm:coverduality} exclude every point outside the cone.
\end{proof}

\subsection{Extension from wire semirings}

Let $|V|\ge2$ and let $\mathcal W_V$ be the subsemiring of the tensor semiring
generated by the binary wires and the unit, under Kronecker
product and \emph{direct sum}, with the restriction preorder
inherited from all tensors. This preorder is Strassen:
natural-number diagonal tensors retain their usual ordering,
every nonzero tensor restricts to the unit, and a tensor of finite
rank restricts from the corresponding diagonal tensor.

\begin{proposition}\label{prop:wiresubsemiring}
Every normalized, additive, multiplicative,
restriction-monotone function on $\mathcal W_V$ extends to an
ordinary spectral point on all $V$-tensors.
The same statement holds when any fixed finite collection of
hub parties is included in $V$.
\end{proposition}

\begin{proof}
The inclusion is a semiring homomorphism and an order embedding
because the preorder was inherited, not generated by a smaller
list of permitted routing moves. Apply the spectral extension
theorem \cite[Theorem 4.16, pp.~53, 56--57 of the author
version]{wz-spectra}. Its hypotheses were checked above.
\end{proof}

Thus assigning a desired cost to each wire is not enough: the
assignment must extend additively to direct sums of wire products
and respect every restriction between those sums.
Once this is proved, extension to all tensors is not a further
obstruction. The finite certificate theorem identifies precisely
what a failure of this compatibility can produce.

\subsection{Why a linear inequality cannot exclude a Euclidean price}

\begin{proposition}[The linear obstruction is insufficient]
\label{prop:no-linear-euclidean-obstruction}
Fix a finite set of labels and any field. Every homogeneous linear
inequality in their pairwise distances that holds for all prices
also holds for every Euclidean configuration of those labels.
In particular a failure of Euclidean flatness cannot be proved by
one homogeneous linear inequality in the lengths.
\end{proposition}

\begin{proof}
For every nontrivial subset $S$ of the labels, the flattening rank
across $S$ is a spectral point. Its wire metric is the cut metric
$\delta_S$, equal to one between $S$ and its complement and zero
otherwise.

Let the labels have positions $x_i\in\R^d$. For a unit vector
$\theta$ and a real number $t$, let
$S_{\theta,t}=\{i:\theta\cdot x_i>t\}$. Then
\[
 \int_{\R}\delta_{S_{\theta,t}}(i,j)\,dt
       =|\theta\cdot(x_i-x_j)|.
\]
Averaging over the unit sphere gives
$\kappa_d\|x_i-x_j\|$, with
$\kappa_d=\int|\theta_1|\,d\theta>0$.
There are only finitely many distinct cuts of a finite label set.
Collecting the measure of parameters producing each cut expresses
the Euclidean metric as a finite nonnegative linear combination
of cut metrics, after division by $\kappa_d$.
The coefficients are finite: nontrivial cuts occur only for offsets
between the smallest and largest projected positions, an interval
bounded by the diameter.
Any homogeneous linear inequality valid on the price cone holds
on each cut metric and hence on their nonnegative sum.
\end{proof}

The maximum over different covering layouts in
Theorem~\ref{thm:finitecertificate} is therefore essential.
Euclidean metrics belong to the convex conic hull of prices;
membership in that hull does not prove membership in the price
cone itself.
